% ARCHIVED v3.80 (2026-09-24): the UAV-corridor block of chapters/06_results.tex as it stood at v3.79 (placeholders pending R33).
% Replaced by the filled block from DA_20260924_corridor_v3.md. Lines 1534-1791 of 06_results.tex before the edit.
\subsection{UAV-corridor}\label{sec:res:uav:corridor}

\subsubsection{Generative Models}\label{sec:res:uav:corridor:raw}
The two constraints are read one at a time, as the avoiding geometries are: every table of this scene has a
\emph{tilt} block and a \emph{hump} block, and nothing is averaged across them.
\hole{Result sentence, before projection: which models cross the finish line unprojected and how many
flights violate each constraint. Read from \autoref{tab:uav-corridor-raw} once the run lands.
The v3.65 reading on the v2 corpus was: every model succeeds on every flight and none is violation-free,
so the scene is a repair case.} The cells that reach the finish line here are the cells whose projected
rows in \autoref{sec:res:uav:corridor} are read; a cost frontier is drawn for this arm only if the models
do not all cross the line, because a frontier over points that share one outcome is a ranking by price
(\autoref{sec:res:aligning:budget}).

\begin{table}[htpb]
  \caption[UAV-corridor: the generative models before projection]{UAV-corridor with the test-time
  constraints, \textbf{before projection}: the plan flown as generated, ten flights per row, spread over
  the three routes, one training seed, one block per constraint (\emph{tilt}, \emph{hump}). The grid is: the two average-velocity models at
  $\nfe\in\{1,2,3\}$, instantaneous-velocity matching at $\nfe\in\{1,2,3,5,20\}$, the diffusion baseline
  at its trained budget, printed last. \emph{Success}: crosses the finish line. \emph{Violation-free}: no
  violating control step. \emph{S\&C}: success on a violation-free flight. \emph{Violating steps} and
  \emph{ms/step}: mean $\pm$ the sample standard deviation over the flights
  (\autoref{sec:setup:metrics:spread}). \emph{Pending}: run not yet landed (R33). \baselinemark\ the
  diffusion model of \ac{DPCC}.}
  \label{tab:uav-corridor-raw}
  \centering
  \footnotesize
  \setlength{\tabcolsep}{3pt}
  \begin{tabularx}{\linewidth}{@{}L r rrr rr@{}}
    \toprule
    Model & $\nfe$ & Success & \shortstack{Violation-\\free} & S\&C & \shortstack{Violating\\steps} & ms/step \\
    \midrule
    \multicolumn{7}{@{}l}{\emph{tilt}}\\
    MeanFM & 1 & \multicolumn{5}{l}{\emph{pending}} \\
    MeanFM & 2 & \multicolumn{5}{l}{\emph{pending}} \\
    MeanFM & 3 & \multicolumn{5}{l}{\emph{pending}} \\
    CI-MeanFM, $\alpha_{\mathrm{end}}=0.2$ & 1 & \multicolumn{5}{l}{\emph{pending}} \\
    CI-MeanFM, $\alpha_{\mathrm{end}}=0.2$ & 2 & \multicolumn{5}{l}{\emph{pending}} \\
    CI-MeanFM, $\alpha_{\mathrm{end}}=0.2$ & 3 & \multicolumn{5}{l}{\emph{pending}} \\
    FM & 1 & \multicolumn{5}{l}{\emph{pending}} \\
    FM & 2 & \multicolumn{5}{l}{\emph{pending}} \\
    FM & 3 & \multicolumn{5}{l}{\emph{pending}} \\
    FM & 5 & \multicolumn{5}{l}{\emph{pending}} \\
    FM & 20 & \multicolumn{5}{l}{\emph{pending}} \\
    \mainconfig{\baselinemark\ Diffusion} & \mainconfig{20} & \multicolumn{5}{l}{\mainconfig{\emph{pending}}} \\
    \midrule
    \multicolumn{7}{@{}l}{\emph{hump}}\\
    MeanFM & 1 & \multicolumn{5}{l}{\emph{pending}} \\
    MeanFM & 2 & \multicolumn{5}{l}{\emph{pending}} \\
    MeanFM & 3 & \multicolumn{5}{l}{\emph{pending}} \\
    CI-MeanFM, $\alpha_{\mathrm{end}}=0.2$ & 1 & \multicolumn{5}{l}{\emph{pending}} \\
    CI-MeanFM, $\alpha_{\mathrm{end}}=0.2$ & 2 & \multicolumn{5}{l}{\emph{pending}} \\
    CI-MeanFM, $\alpha_{\mathrm{end}}=0.2$ & 3 & \multicolumn{5}{l}{\emph{pending}} \\
    FM & 1 & \multicolumn{5}{l}{\emph{pending}} \\
    FM & 2 & \multicolumn{5}{l}{\emph{pending}} \\
    FM & 3 & \multicolumn{5}{l}{\emph{pending}} \\
    FM & 5 & \multicolumn{5}{l}{\emph{pending}} \\
    FM & 20 & \multicolumn{5}{l}{\emph{pending}} \\
    \mainconfig{\baselinemark\ Diffusion} & \mainconfig{20} & \multicolumn{5}{l}{\mainconfig{\emph{pending}}} \\
    \bottomrule
  \end{tabularx}
\end{table}
\dataref{corridor runs, R33 (\texttt{PENDING\_20260923\_uav\_corridor\_v3\_pillars\_v2\_paper\_runs.md}), geometries \texttt{corridor\_v3\_tilt} and \texttt{corridor\_v3\_ablation\_hump}, variant
\texttt{diffuser}: \texttt{success\_relaxed}, \texttt{collision\_free\_completed},
\texttt{n\_success\_relaxed\_and\_constraints}, \texttt{n\_violations}, \texttt{avg\_time\_ms}. The v2
corpus this table was computed from until v3.65 is archived.}

\FloatBarrier

\FloatBarrier
\subsubsection{Projection}\label{sec:res:uav:projected}
\hole{Result sentence, after projection: which models succeed violation-free at which budget under the
per-step projector of \ac{DPCC} at activation threshold $0.5$, and where endpoint projection, which has a
guiding step from $\nfe=3$, stands beside it (\autoref{tab:uav-corridor}). Read the two projectors
together and pick the best configuration of each model by the table and the frontier.} The budgets are
read as one ladder across the scene: the per-step projector is active at every budget, endpoint projection
at $\nfe\ge 3$, where it has a guiding step (\autoref{sec:res:constraints:degenerate}), so at $\nfe=1$ and
$\nfe=2$ the table carries per-step rows only.

\begin{table}[htpb]
  \caption[UAV-corridor: the generative models after projection]{UAV-corridor with its two test-time
  constraints, one block each, \textbf{after projection}: each model at each budget of the grid under the per-step
  projector of \ac{DPCC} (activation threshold $0.5$, tightened constraints) and, from $\nfe=3$, under
  endpoint projection, each at its best selection rule of \autoref{tab:uav-corridor-projection}; ten
  flights per cell, one training seed. \emph{S\&C}: success on a violation-free flight; \emph{ms/step}:
  the time to compute one action. \emph{---}: endpoint projection has no guiding step at that budget.
  \emph{Pending}: run not yet landed (R33). \baselinemark\ the diffusion model of \ac{DPCC}.}
  \label{tab:uav-corridor}
  \centering
  \footnotesize
  \setlength{\tabcolsep}{4pt}
  \begin{tabularx}{\linewidth}{@{}L r rrr rrr@{}}
    \toprule
    & & \multicolumn{3}{c}{per-step projection} & \multicolumn{3}{c}{endpoint projection} \\
    \cmidrule(lr){3-5}\cmidrule(l){6-8}
    Model & $\nfe$ & S\&C & rule & ms/step & S\&C & rule & ms/step \\
    \midrule
    \multicolumn{8}{@{}l}{\emph{tilt}}\\
    MeanFM & 1 & \multicolumn{3}{l}{\emph{pending}} & --- & --- & --- \\
    MeanFM & 2 & \multicolumn{3}{l}{\emph{pending}} & --- & --- & --- \\
    MeanFM & 3 & \multicolumn{3}{l}{\emph{pending}} & \multicolumn{3}{l}{\emph{pending}} \\
    CI-MeanFM, $\alpha_{\mathrm{end}}=0.2$ & 1 & \multicolumn{3}{l}{\emph{pending}} & --- & --- & --- \\
    CI-MeanFM, $\alpha_{\mathrm{end}}=0.2$ & 2 & \multicolumn{3}{l}{\emph{pending}} & --- & --- & --- \\
    CI-MeanFM, $\alpha_{\mathrm{end}}=0.2$ & 3 & \multicolumn{3}{l}{\emph{pending}} & \multicolumn{3}{l}{\emph{pending}} \\
    FM & 1 & \multicolumn{3}{l}{\emph{pending}} & --- & --- & --- \\
    FM & 2 & \multicolumn{3}{l}{\emph{pending}} & --- & --- & --- \\
    FM & 3 & \multicolumn{3}{l}{\emph{pending}} & \multicolumn{3}{l}{\emph{pending}} \\
    FM & 5 & \multicolumn{3}{l}{\emph{pending}} & \multicolumn{3}{l}{\emph{pending}} \\
    FM & 20 & \multicolumn{3}{l}{\emph{pending}} & \multicolumn{3}{l}{\emph{pending}} \\
    \mainconfig{\baselinemark\ Diffusion} & \mainconfig{20} & \multicolumn{3}{l}{\mainconfig{\emph{pending}}} & --- & --- & --- \\
    \midrule
    \multicolumn{8}{@{}l}{\emph{hump}}\\
    MeanFM & 1 & \multicolumn{3}{l}{\emph{pending}} & --- & --- & --- \\
    MeanFM & 2 & \multicolumn{3}{l}{\emph{pending}} & --- & --- & --- \\
    MeanFM & 3 & \multicolumn{3}{l}{\emph{pending}} & \multicolumn{3}{l}{\emph{pending}} \\
    CI-MeanFM, $\alpha_{\mathrm{end}}=0.2$ & 1 & \multicolumn{3}{l}{\emph{pending}} & --- & --- & --- \\
    CI-MeanFM, $\alpha_{\mathrm{end}}=0.2$ & 2 & \multicolumn{3}{l}{\emph{pending}} & --- & --- & --- \\
    CI-MeanFM, $\alpha_{\mathrm{end}}=0.2$ & 3 & \multicolumn{3}{l}{\emph{pending}} & \multicolumn{3}{l}{\emph{pending}} \\
    FM & 1 & \multicolumn{3}{l}{\emph{pending}} & --- & --- & --- \\
    FM & 2 & \multicolumn{3}{l}{\emph{pending}} & --- & --- & --- \\
    FM & 3 & \multicolumn{3}{l}{\emph{pending}} & \multicolumn{3}{l}{\emph{pending}} \\
    FM & 5 & \multicolumn{3}{l}{\emph{pending}} & \multicolumn{3}{l}{\emph{pending}} \\
    FM & 20 & \multicolumn{3}{l}{\emph{pending}} & \multicolumn{3}{l}{\emph{pending}} \\
    \mainconfig{\baselinemark\ Diffusion} & \mainconfig{20} & \multicolumn{3}{l}{\mainconfig{\emph{pending}}} & --- & --- & --- \\
    \bottomrule
  \end{tabularx}
\end{table}
\dataref{corridor runs, R33, both geometries; per-step = \texttt{dpcc-\{r,c,t\}} at $\eta=0.5$, endpoint =
\texttt{hardflow\_sls\{,-r,-c,-t\}}, both on the corridor's projection stack
(\autoref{sec:setup:protocol:uav}); the best rule per cell is the one \autoref{tab:uav-corridor-projection}
prints highest. The unprojected twins are \autoref{tab:uav-corridor-raw}.}

\hole{Budget paragraph: at which budget each model stops violating the halfspace, whether that is the same
budget at which it stops reaching the finish line, and what the diffusion baseline does at its one budget.
Then the frontier reading of \autoref{fig:uav-corridor-tradeoff}: which configurations are on it, the
cheapest one, and whether its shape is a step or a slope.} The frontier is drawn as before: success with
constraint satisfaction against the time one action costs, logarithmic in the number of successful flights,
a configuration with no successful flight not eligible, costs within one per cent treated as one; both
projectors are drawn on it, per-step as circles and endpoint as squares, so that the choice of projector is
read off the same picture as the choice of budget.

\begin{figure}[htpb]
  \centering
  \todofigure[0.34\textwidth]{\texttt{fig\_uav\_corridor\_tradeoff}, one panel per constraint (tilt, hump): success with
  constraint satisfaction (log axis in successful flights of ten, floor for zero) against the time to
  compute one action, every cell of \autoref{tab:uav-corridor-projection} under both projectors, one
  training seed; rings and staircase on the frontier; builder
  \texttt{DA\_in\_Paper/plotting/builders/frontier.py::fig\_uav\_corridor\_tradeoff} extended to the
  endpoint cells. Until R33 lands the builder draws the v2 corpus, which is archived and not this figure.}
  \caption[UAV-corridor: outcome against cost]{UAV-corridor, success with constraint satisfaction against
  the time to compute one action, both projectors, every budget of the grid, one training seed; drawn as
  in \autoref{fig:avoiding-tradeoff}. Pending R33.}
  \label{fig:uav-corridor-tradeoff}
\end{figure}

\hole{Paths paragraph: what the flown paths of \autoref{fig:uav-corridor-paths} show --- the budget at
which a model stops crossing the halfspace against the budget at which it stops reaching the end of the
corridor, and where the diffusion baseline sits on that trade.}

\begin{figure}[htpb]
  \centering
  \todofigure[0.34\textwidth]{\texttt{fig\_uav\_corridor\_paths}, one row per constraint (tilt, hump): the scene from above with
  the ten flown paths per cell, one row per model across its budgets and one panel per projector at the
  best rule; colour = constraint, stroke and end mark = success, cross = collided or lost control; extract
  \texttt{plotting/extract/uav\_paths.py} (group \texttt{corridor}) on the staged rollouts of R33.}
  \caption[UAV-corridor: flown paths per model and budget]{UAV-corridor seen from above with the executed
  path of each flight, per model, budget and projector. Pending R33.}
  \label{fig:uav-corridor-paths}
\end{figure}

\paragraph{Altitude.} Both constraints of this scene act on altitude, and that is what they are for: under
the tilt the room is below the launch altitude, so a projected flight has to descend as it moves sideways;
under the hump it has to climb. \autoref{fig:uav-corridor-altitude} shows the flights from the side so that
this is seen rather than assumed. \hole{Altitude numbers per constraint: the flown band, the paired
difference between the projected and the unprojected flight, and where along $x$ it occurs. In the pilots
the projected flights descended $0.15$--$0.29$\,m under the tilt and climbed $0.30$--$0.40$\,m over the hump,
where on the level corridor of the earlier evaluation the projection moved altitude only as a by-product
of the lateral correction.}

\begin{figure}[htpb]
  \centering
  \todofigure[0.30\textwidth]{\texttt{fig\_uav\_corridor\_altitude}, one row per constraint (tilt: the descent under the leaned plane; hump: the climb): altitude $z$ over $x$
  for every flight, one panel per model, grey unprojected and blue projected, the inflated box limits
  dashed; builder \texttt{plotting/builders/altitude.py} on the R33 rollouts (group
  \texttt{corridor\_altitude}).}
  \caption[UAV-corridor: flown altitude, without and with projection]{UAV-corridor seen from the side: the
  altitude of every flight over its position along the corridor, without and with projection, one panel
  per model. Pending R33.}
  \label{fig:uav-corridor-altitude}
\end{figure}

UAV-pillars is not drawn as a frontier: its result is the difference between a plan's outcome on
D3IL-avoiding and the outcome of the same plan flown, which is a comparison and not a trade-off
(\autoref{sec:res:uav:models}). UAV-s-curve has data, but every
configuration it has --- projected and unprojected, at every budget flown --- scores $0.00$ success with
constraint satisfaction (\autoref{tab:uav-scurve}). A cost frontier
over points that all share one outcome is a ranking of failures by price, so what that scene gets instead
is the abort count printed beside each of its rows.

\FloatBarrier

\FloatBarrier

\FloatBarrier
\subsubsection{Projection Methods}\label{sec:res:uav:projection}
\hole{Corridor result sentence, per-step against endpoint: which projector is the method on this scene,
under which rule and at which budget, read from \autoref{tab:uav-corridor-projection}; whether endpoint
projection clears the bar this thesis set for it --- overtaking the per-step projector of \ac{DPCC} at a
lower projection cost --- and whether either projector changes with the budget.}

\begin{table}[htpb]
  \caption[UAV-corridor: the two projection methods under every selection rule]{UAV-corridor, success with
  constraint satisfaction --- success on a violation-free flight --- for both projection methods under
  every selection rule, one block per constraint (\emph{tilt}, \emph{hump}), at the budgets at which
  endpoint projection has a guiding step; ten flights
  per cell, tightened constraints, activation threshold $0.5$, one training seed. Per-step: $r$ random,
  $c$ cumulative projection cost, $t$ temporal consistency, four candidate plans. Endpoint:
  \emph{single} one candidate, then the three rules over four. The diffusion baseline has per-step rows
  only. \emph{Pending}: R33.}
  \label{tab:uav-corridor-projection}
  \centering
  \footnotesize
  \setlength{\tabcolsep}{5pt}
  \begin{tabularx}{\linewidth}{@{}L r rrr rrrr@{}}
    \toprule
    & & \multicolumn{3}{c}{Per-step} & \multicolumn{4}{c}{Endpoint} \\
    \cmidrule(lr){3-5}\cmidrule(l){6-9}
    Model & $\nfe$ & $r$ & $c$ & $t$ & single & $r$ & $c$ & $t$ \\
    \midrule
    \multicolumn{9}{@{}l}{\emph{tilt}}\\
    MeanFM & 3 & \multicolumn{7}{l}{\emph{pending}} \\
    CI-MeanFM, $\alpha_{\mathrm{end}}=0.2$ & 3 & \multicolumn{7}{l}{\emph{pending}} \\
    FM & 3 & \multicolumn{7}{l}{\emph{pending}} \\
       & 5 & \multicolumn{7}{l}{\emph{pending}} \\
       & 20 & \multicolumn{7}{l}{\emph{pending}} \\
    \mainconfig{\baselinemark\ Diffusion} & \mainconfig{20} & \multicolumn{3}{l}{\mainconfig{\emph{pending}}} & --- & --- & --- & --- \\
    \midrule
    \multicolumn{9}{@{}l}{\emph{hump}}\\
    MeanFM & 3 & \multicolumn{7}{l}{\emph{pending}} \\
    CI-MeanFM, $\alpha_{\mathrm{end}}=0.2$ & 3 & \multicolumn{7}{l}{\emph{pending}} \\
    FM & 3 & \multicolumn{7}{l}{\emph{pending}} \\
       & 5 & \multicolumn{7}{l}{\emph{pending}} \\
       & 20 & \multicolumn{7}{l}{\emph{pending}} \\
    \mainconfig{\baselinemark\ Diffusion} & \mainconfig{20} & \multicolumn{3}{l}{\mainconfig{\emph{pending}}} & --- & --- & --- & --- \\
    \bottomrule
  \end{tabularx}
\end{table}
\dataref{corridor runs, R33, both geometries: variants \texttt{dpcc-\{r,c,t\}} and \texttt{hardflow\_sls\{,-r,-c,-t\}}
on the corridor projection stack; \texttt{n\_success\_relaxed\_and\_constraints}.}
\guard{Endpoint projection has no row at $\nfe=1$ or $\nfe=2$ on this scene and will not get one: at the
activation threshold of $0.5$ the only active sampling step there is the terminal one, so there is no
guiding step for the method to act on (\autoref{sec:res:constraints:degenerate}). That is a property of
the method at those budgets, not a pending evaluation. The baseline's dashes are the opposite case:
endpoint projection has never been run on the diffusion model on any scene.}

The projection configuration of this scene differs from that of the other quadrotor scene
(\autoref{sec:setup:protocol:uav}), so the scenes are not pooled.

\FloatBarrier
